% !TEX encoding = UTF-8
%======================================================================
% 1. LỚP TÀI LIỆU VÀ ENCODING
%======================================================================
\documentclass[12pt, a4paper, oneside]{report}

% Hỗ trợ gõ tiếng Việt và hiển thị đúng
\usepackage[utf8]{vietnam}
\usepackage[T5]{fontenc}
\usepackage{textcomp}

% Thiết lập font chữ Times New Roman hoặc tương đương
\usepackage{mathptmx} 

%======================================================================
% 2. CANH LỀ VÀ LAYOUT TRANG
%======================================================================
\usepackage[top=3cm, bottom=3cm, left=3.5cm, right=2cm, heightrounded, headheight=15pt]{geometry}

\usepackage{fancyhdr}
\pagestyle{fancy}
\fancyhf{} 
\lhead{\footnotesize\leftmark} 
\cfoot{\footnotesize\thepage} 
\renewcommand{\headrulewidth}{0.4pt}
\renewcommand{\footrulewidth}{0.4pt}

\usepackage{setspace}
\onehalfspacing
\usepackage{indentfirst}

%======================================================================
% 3. CÁC GÓI TOÁN HỌC
%======================================================================
\usepackage{amsmath}    
\usepackage{amssymb}    
\usepackage{amsthm}     
\usepackage{latexsym}   
\usepackage{amsbsy}     
\numberwithin{equation}{chapter} 

\newtheorem{definition}{Định nghĩa}[chapter]
\newtheorem{theorem}{Định lý}[chapter]
\newtheorem{corollary}{Hệ quả}[chapter]
\newtheorem{lemma}{Bổ đề}[chapter]

%======================================================================
% 4. HÌNH ẢNH, MÀU SẮC & GIAO DIỆN
%======================================================================
\usepackage{graphicx}   
\usepackage{svg}        
\usepackage[table, svgnames, xcdraw]{xcolor} % Gộp tất cả option màu vào 1 dòng
\usepackage{tikz}       
\graphicspath{{./images/}} 

\usepackage[labelfont=bf, font=small, skip=5pt]{caption}
\usepackage{subcaption} 
\usepackage{float}      % Hỗ trợ lệnh [H] cho ảnh/bảng
\usepackage{framed}     % [RẤT QUAN TRỌNG] Hỗ trợ môi trường leftbar
\usepackage[most]{tcolorbox}
\usepackage{pdfpages}

\counterwithin{figure}{chapter} 
\counterwithin{table}{chapter}  

%======================================================================
% 5. BẢNG BIỂU
%======================================================================
\usepackage{booktabs}   
\usepackage{tabularx}   
\usepackage{multirow}   
\usepackage{longtable}  
\usepackage{array}      

%======================================================================
% 6. ĐỊNH DẠNG CODE & LIST
%======================================================================
\usepackage{verbatim}
\usepackage{enumitem}

% Sử dụng minted (Lưu ý: Yêu cầu -shell-escape khi biên dịch)
\usepackage{minted}
\setminted{
    frame=lines,      
    framesep=2mm,     
    baselinestretch=1.2,
    fontsize=\small,  
    linenos           
}

% Định dạng lại tiêu đề chương, mục
\usepackage{titlesec}
\titleformat{\chapter}[display]
  {\normalfont\huge\bfseries\centering}{\chaptertitlename\ \thechapter}{20pt}{\Huge}
\titlespacing*{\chapter}{0pt}{50pt}{40pt}

%======================================================================
% 7. TÀI LIỆU THAM KHẢO VÀ HYPERLINK (PHẢI NẰM CUỐI CÙNG)
%======================================================================
\usepackage{csquotes}
\usepackage[
    backend=biber,      
    style=numeric,      
    sorting=none,       
    giveninits=true     
]{biblatex}
\addbibresource{bibliography.bib}

\usepackage{url}        
\usepackage[
    colorlinks=true,      
    linkcolor=black,      
    citecolor=black,      
    urlcolor=black,       
    bookmarks=true,
    pdfencoding=auto,
    unicode
]{hyperref}